% ──────────────────────────────────────────────────────────────
\subsection{Slicing by an arbitrary spacelike hypersurface}%
\label{sec:general-slicing}

The equivalence of the covariant and $3{+}1$ forms does not depend
on the particular choice of time coordinate~$x^0$.  We now state
the general result, which applies to arbitrary spacelike
hypersurfaces and extends naturally to curved spacetimes.

\begin{theorem}[General slicing theorem]%
\label{thm:general-slicing}
  Let $\Sigma = f^{-1}(0)$ be a smooth spacelike Cauchy surface
  in~$(\Rn{4},\eta)$, defined as a level set of a smooth function
  $f\colon\Rn{4}\to\R$ with $\nabla f$ everywhere timelike
  (that is, $\eta^{\alpha\beta}\partial_\alpha f\,\partial_\beta f < 0$).
  Then each worldline~$\gamma_j$ intersects~$\Sigma$
  transversally at a unique point
  $\gamma_j(\tau_j^\ast)$, and the restriction (``slice'') of
  $T^{\alpha\beta}$ to~$\Sigma$ satisfies
  \begin{equation}\label{eq:general-slicing}
    T^{\alpha\beta}\big|_\Sigma(\vb{y})
      = \sum_{j=1}^{N}
        \frac{p_j^\alpha(\tau_j^\ast)\,
              \dot\gamma_j^\beta(\tau_j^\ast)}
             {\bigl|\partial_\gamma f\,
              \dot\gamma_j^\gamma(\tau_j^\ast)\bigr|}\,
        \delta_\Sigma(\vb{y} - \vb{y}_j)\,,
  \end{equation}
  where $\vb{y}$ denotes coordinates on~$\Sigma$, $\vb{y}_j$ is
  the location of the $j$-th crossing on~$\Sigma$, and
  $\delta_\Sigma$ is the delta function with respect to the
  Riemannian measure on~$\Sigma$ induced by the ambient Minkowski
  metric; since $\Sigma$ is spacelike, this induced metric is
  positive definite.  (This coincides with the $3$-dimensional
  Hausdorff measure~$\mathcal{H}^3$ restricted to~$\Sigma$, up
  to identification of~$\Sigma$ with~$\R^3$ via coordinates.)
\end{theorem}

\begin{proof}[Sketch of proof]
  \emph{Transversality.}\enspace
  The covector field $\partial_\alpha f$ satisfies
  $\eta^{\alpha\beta}\partial_\alpha f\,\partial_\beta f < 0$
  (timelike), and the tangent vector $\dot\gamma^\alpha$ is also
  timelike.  By the reverse Cauchy--Schwarz inequality for
  Lorentzian inner products (see, e.g.,
  O'Neill~\cite[Prop.~5.30]{ONeill83}: if $u,v$ are both
  timelike and in the same time-cone, decompose in an orthonormal
  frame where $u = (u^0,0,0,0)$; then
  $|g(u,v)| = |u^0||v^0| \ge |u||v|$ since
  $|v^0| \ge |v_{\text{spatial}}|$), the contraction
  $\partial_\gamma f \cdot \dot\gamma^\gamma \neq 0$.  This
  establishes transversality.

  \emph{Uniqueness of intersection.}\enspace
  Consider the composition $h = f\circ\gamma\colon\R\to\R$.
  Its derivative $h'(\tau) = \partial_\gamma f\bigl(\gamma(\tau)\bigr)
  \,\dot\gamma^\gamma(\tau)$ is nonzero everywhere by
  transversality, and has constant sign by continuity (since
  $\R$ is connected).  Hence $h$ is strictly monotone, so
  $h^{-1}(0)$ contains at most one point.  Existence of the
  crossing follows from the Cauchy-surface hypothesis: by
  definition, every inextendible timelike curve meets a Cauchy
  surface, so each complete worldline~$\gamma_j$ intersects
  $\Sigma = f^{-1}(0)$ at least---hence exactly---once.

  \emph{Slicing and the coarea formula.}\enspace
  The proof proceeds in two stages.  First, the coarea formula
  applied to $h = f\circ\gamma\colon\R\to\R$---which in one
  dimension is simply the classical change-of-variables
  (substitution) formula---converts the proper-time integral into
  an integral with respect to the level-set parameter of~$f$.
  Second, the identification of the resulting object as a delta
  function on~$\Sigma$ uses the theory of slicing of normal
  $1$-currents (Federer~\cite[Theorem~4.3.2]{Federer69}), which
  provides the natural framework when $\Sigma$ is not a
  coordinate hyperplane.  The Jacobian factor
  $|\partial_\gamma f \cdot \dot\gamma^\gamma|^{-1}$ arises from
  the first stage.
\end{proof}

\begin{remark}
  For the standard time function $f(x) = x^0 - t_0$, we have
  $\partial_\gamma f = \delta_\gamma^0$, so
  $|\partial_\gamma f\cdot\dot\gamma^\gamma| = |\dot\gamma^0|$.
  Since $\gamma$ is future-directed
  (Definition~\ref{def:timelike-worldline}),
  $|\dot\gamma^0| = d\gamma^0/d\tau$, and the Jacobian factor
  combines with $\dot\gamma^\beta$ to give
  $d\gamma^\beta/dt$---recovering
  Theorem~\ref{thm:equivalence}.  In curved spacetime one
  replaces $f$ with a Cauchy time function and $\delta_\Sigma$
  with the delta function defined by the induced Riemannian metric
  on the Cauchy surface.
\end{remark}
